%!TeX root=../gardentop.tex
\chapter{Dickon} 
\begin{figure}[t!]
\centering
\includegraphics[width=\linewidth]{september}
\end{figure}

 \lettrine[lines=6]{T}{he} sun shone down for nearly a week on the secret garden. The Secret Garden was what Mary called it when she was thinking of it. She liked the name, and she liked still more the feeling that when its beautiful old walls shut her in no one knew where she was. It seemed almost like being shut out of the world in some fairy place. The few books she had read and liked had been fairy-story books, and she had read of secret gardens in some of the stories. Sometimes people went to sleep in them for a hundred years, which she had thought must be rather stupid. She had no intention of going to sleep, and, in fact, she was becoming wider awake every day which passed at Misselthwaite. She was beginning to like to be out of doors; she no longer hated the wind, but enjoyed it. She could run faster, and longer, and she could skip up to a hundred. The bulbs in the secret garden must have been much astonished. Such nice clear places were made round them that they had all the breathing space they wanted, and really, if Mistress Mary had known it, they began to cheer up under the dark earth and work tremendously. The sun could get at them and warm them, and when the rain came down it could reach them at once, so they began to feel very much alive.

Mary was an odd, determined little person, and now she had something interesting to be determined about, she was very much absorbed, indeed. She worked and dug and pulled up weeds steadily, only becoming more pleased with her work every hour instead of tiring of it. It seemed to her like a fascinating sort of play. She found many more of the sprouting pale green points than she had ever hoped to find. They seemed to be starting up everywhere and each day she was sure she found tiny new ones, some so tiny that they barely peeped above the earth. There were so many that she remembered what Martha had said about the »snowdrops by the thousands,« and about bulbs spreading and making new ones. These had been left to themselves for ten years and perhaps they had spread, like the snowdrops, into thousands. She wondered how long it would be before they showed that they were flowers. Sometimes she stopped digging to look at the garden and try to imagine what it would be like when it was covered with thousands of lovely things in bloom.

During that week of sunshine, she became more intimate with Ben Weatherstaff. She surprised him several times by seeming to start up beside him as if she sprang out of the earth. The truth was that she was afraid that he would pick up his tools and go away if he saw her coming, so she always walked toward him as silently as possible. But, in fact, he did not object to her as strongly as he had at first. Perhaps he was secretly rather flattered by her evident desire for his elderly company. Then, also, she was more civil than she had been. He did not know that when she first saw him she spoke to him as she would have spoken to a native, and had not known that a cross, sturdy old Yorkshire man was not accustomed to salaam to his masters, and be merely commanded by them to do things.

»Tha'rt like th' robin,« he said to her one morning when he lifted his head and saw her standing by him. »I never knows when I shall see thee or which side tha'll come from.«

»He's friends with me now,« said Mary.

»That's like him,« snapped Ben Weatherstaff. »Makin' up to th' women folk just for vanity an' flightiness. There's nothin' he wouldn't do for th' sake o' showin' off an' flirtin' his tail-feathers. He's as full o' pride as an egg's full o' meat.«

He very seldom talked much and sometimes did not even answer Mary's questions except by a grunt, but this morning he said more than usual. He stood up and rested one hobnailed boot on the top of his spade while he looked her over.

»How long has tha' been here?« he jerked out.

»I think it's about a month,« she answered.

»Tha's beginnin' to do Misselthwaite credit,« he said. »Tha's a bit fatter than tha' was an' tha's not quite so yeller. Tha' looked like a young plucked crow when tha' first came into this garden. Thinks I to myself I never set eyes on an uglier, sourer faced young 'un.«

Mary was not vain and as she had never thought much of her looks she was not greatly disturbed.

»I know I'm fatter,« she said. »My stockings are getting tighter. They used to make wrinkles. There's the robin, Ben Weatherstaff.«

There, indeed, was the robin, and she thought he looked nicer than ever. His red waistcoat was as glossy as satin and he flirted his wings and tail and tilted his head and hopped about with all sorts of lively graces. He seemed determined to make Ben Weatherstaff admire him. But Ben was sarcastic.

»Aye, there tha' art!« he said. »Tha' can put up with me for a bit sometimes when tha's got no one better. Tha's been reddinin' up thy waistcoat an' polishin' thy feathers this two weeks. I know what tha's up to. Tha's courtin' some bold young madam somewhere, tellin' thy lies to her about bein' th' finest cock robin on Missel Moor an' ready to fight all th' rest of 'em.«

»Oh! look at him!« exclaimed Mary.

The robin was evidently in a fascinating, bold mood. He hopped closer and closer and looked at Ben Weatherstaff more and more engagingly. He flew on to the nearest currant bush and tilted his head and sang a little song right at him.

»Tha' thinks tha'll get over me by doin' that,« said Ben, wrinkling his face up in such a way that Mary felt sure he was trying not to look pleased. »Tha' thinks no one can stand out against thee—that's what tha' thinks.«

The robin spread his wings—Mary could scarcely believe her eyes. He flew right up to the handle of Ben Weatherstaff's spade and alighted on the top of it. Then the old man's face wrinkled itself slowly into a new expression. He stood still as if he were afraid to breathe—as if he would not have stirred for the world, lest his robin should start away. He spoke quite in a whisper.

»Well, I'm danged!« he said as softly as if he were saying something quite different. »Tha' does know how to get at a chap—tha' does! Tha's fair unearthly, tha's so knowin'.«

And he stood without stirring—almost without drawing his breath—until the robin gave another flirt to his wings and flew away. Then he stood looking at the handle of the spade as if there might be Magic in it, and then he began to dig again and said nothing for several minutes.

But because he kept breaking into a slow grin now and then, Mary was not afraid to talk to him.

»Have you a garden of your own?« she asked.

»No. I'm bachelder an' lodge with Martin at th' gate.«

»If you had one,« said Mary, »what would you plant?«

»Cabbages an' 'taters an' onions.«

»But if you wanted to make a flower garden,« persisted Mary, »what would you plant?«

»Bulbs an' sweet-smellin' things—but mostly roses.«

Mary's face lighted up.

»Do you like roses?« she said.

Ben Weatherstaff rooted up a weed and threw it aside before he answered.

»Well, yes, I do. I was learned that by a young lady I was gardener to. She had a lot in a place she was fond of, an' she loved 'em like they was children—or robins. I've seen her bend over an' kiss 'em.« He dragged out another weed and scowled at it. »That were as much as ten year' ago.«

»Where is she now?« asked Mary, much interested.

»Heaven,« he answered, and drove his spade deep into the soil, »'cording to what parson says.«

»What happened to the roses?« Mary asked again, more interested than ever.

»They was left to themselves.«

Mary was becoming quite excited.

»Did they quite die? Do roses quite die when they are left to themselves?« she ventured.

»Well, I'd got to like 'em—an' I liked her—an' she liked 'em,« Ben Weatherstaff admitted reluctantly. »Once or twice a year I'd go an' work at 'em a bit—prune 'em an' dig about th' roots. They run wild, but they was in rich soil, so some of 'em lived.«

»When they have no leaves and look gray and brown and dry, how can you tell whether they are dead or alive?« inquired Mary.

»Wait till th' spring gets at 'em—wait till th' sun shines on th' rain an' th' rain falls on th' sunshine an' then tha'll find out.«

»How—how?« cried Mary, forgetting to be careful.

»Look along th' twigs an' branches an' if tha' sees a bit of a brown lump swelling here an' there, watch it after th' warm rain an' see what happens.« He stopped suddenly and looked curiously at her eager face. »Why does tha' care so much about roses an' such, all of a sudden?« he demanded.

Mistress Mary felt her face grow red. She was almost afraid to answer.

»I—I want to play that—that I have a garden of my own,« she stammered. »I—there is nothing for me to do. I have nothing—and no one.«

»Well,« said Ben Weatherstaff slowly, as he watched her, »that's true. Tha' hasn't.«

He said it in such an odd way that Mary wondered if he was actually a little sorry for her. She had never felt sorry for herself; she had only felt tired and cross, because she disliked people and things so much. But now the world seemed to be changing and getting nicer. If no one found out about the secret garden, she should enjoy herself always.

She stayed with him for ten or fifteen minutes longer and asked him as many questions as she dared. He answered every one of them in his queer grunting way and he did not seem really cross and did not pick up his spade and leave her. He said something about roses just as she was going away and it reminded her of the ones he had said he had been fond of.

»Do you go and see those other roses now?« she asked.

»Not been this year. My rheumatics has made me too stiff in th' joints.«

He said it in his grumbling voice, and then quite suddenly he seemed to get angry with her, though she did not see why he should.

»Now look here!« he said sharply. »Don't tha' ask so many questions. Tha'rt th' worst wench for askin' questions I've ever come across. Get thee gone an' play thee. I've done talkin' for to-day.«

And he said it so crossly that she knew there was not the least use in staying another minute. She went skipping slowly down the outside walk, thinking him over and saying to herself that, queer as it was, here was another person whom she liked in spite of his crossness. She liked old Ben Weatherstaff. Yes, she did like him. She always wanted to try to make him talk to her. Also she began to believe that he knew everything in the world about flowers.

There was a laurel-hedged walk which curved round the secret garden and ended at a gate which opened into a wood, in the park. She thought she would skip round this walk and look into the wood and see if there were any rabbits hopping about. She enjoyed the skipping very much and when she reached the little gate she opened it and went through because she heard a low, peculiar whistling sound and wanted to find out what it was.

It was a very strange thing indeed. She quite caught her breath as she stopped to look at it. A boy was sitting under a tree, with his back against it, playing on a rough wooden pipe. He was a funny looking boy about twelve. He looked very clean and his nose turned up and his cheeks were as red as poppies and never had Mistress Mary seen such round and such blue eyes in any boy's face. And on the trunk of the tree he leaned against, a brown squirrel was clinging and watching him, and from behind a bush nearby a cock pheasant was delicately stretching his neck to peep out, and quite near him were two rabbits sitting up and sniffing with tremulous noses—and actually it appeared as if they were all drawing near to watch him and listen to the strange low little call his pipe seemed to make.

When he saw Mary he held up his hand and spoke to her in a voice almost as low as and rather like his piping.

»Don't tha' move,« he said. »It'd flight 'em.«

Mary remained motionless. He stopped playing his pipe and began to rise from the ground. He moved so slowly that it scarcely seemed as though he were moving at all, but at last he stood on his feet and then the squirrel scampered back up into the branches of his tree, the pheasant withdrew his head and the rabbits dropped on all fours and began to hop away, though not at all as if they were frightened.

»I'm Dickon,« the boy said. »I know tha'rt Miss Mary.«

Then Mary realized that somehow she had known at first that he was Dickon. Who else could have been charming rabbits and pheasants as the natives charm snakes in India? He had a wide, red, curving mouth and his smile spread all over his face.

»I got up slow,« he explained, »because if tha' makes a quick move it startles 'em. A body 'as to move gentle an' speak low when wild things is about.«

He did not speak to her as if they had never seen each other before but as if he knew her quite well. Mary knew nothing about boys and she spoke to him a little stiffly because she felt rather shy.

»Did you get Martha's letter?« she asked.

He nodded his curly, rust-coloured head.

»That's why I come.«

He stooped to pick up something which had been lying on the ground beside him when he piped.

»I've got th' garden tools. There's a little spade an' rake an' a fork an' hoe. Eh! they are good 'uns. There's a trowel, too. An' th' woman in th' shop threw in a packet o' white poppy an' one o' blue larkspur when I bought th' other seeds.«

»Will you show the seeds to me?« Mary said.

She wished she could talk as he did. His speech was so quick and easy. It sounded as if he liked her and was not the least afraid she would not like him, though he was only a common moor boy, in patched clothes and with a funny face and a rough, rusty-red head. As she came closer to him she noticed that there was a clean fresh scent of heather and grass and leaves about him, almost as if he were made of them. She liked it very much and when she looked into his funny face with the red cheeks and round blue eyes she forgot that she had felt shy.

»Let us sit down on this log and look at them,« she said.

They sat down and he took a clumsy little brown paper package out of his coat pocket. He untied the string and inside there were ever so many neater and smaller packages with a picture of a flower on each one.

»There's a lot o' mignonette an' poppies,« he said. »Mignonette's th' sweetest smellin' thing as grows, an' it'll grow wherever you cast it, same as poppies will. Them as'll come up an' bloom if you just whistle to 'em, them's th' nicest of all.«

He stopped and turned his head quickly, his poppy-cheeked face lighting up.

»Where's that robin as is callin' us?« he said.

The chirp came from a thick holly bush, bright with scarlet berries, and Mary thought she knew whose it was.

»Is it really calling us?« she asked.

»Aye,« said Dickon, as if it was the most natural thing in the world, »he's callin' some one he's friends with. That's same as sayin' 'Here I am. Look at me. I wants a bit of a chat.' There he is in the bush. Whose is he?«

»He's Ben Weatherstaff's, but I think he knows me a little,« answered Mary.

»Aye, he knows thee,« said Dickon in his low voice again. »An' he likes thee. He's took thee on. He'll tell me all about thee in a minute.«

He moved quite close to the bush with the slow movement Mary had noticed before, and then he made a sound almost like the robin's own twitter. The robin listened a few seconds, intently, and then answered quite as if he were replying to a question.

»Aye, he's a friend o' yours,« chuckled Dickon.

»Do you think he is?« cried Mary eagerly. She did so want to know. »Do you think he really likes me?«

»He wouldn't come near thee if he didn't,« answered Dickon. »Birds is rare choosers an' a robin can flout a body worse than a man. See, he's making up to thee now. 'Cannot tha' see a chap?' he's sayin'.«

And it really seemed as if it must be true. He so sidled and twittered and tilted as he hopped on his bush.

»Do you understand everything birds say?« said Mary.

Dickon's grin spread until he seemed all wide, red, curving mouth, and he rubbed his rough head.

»I think I do, and they think I do,« he said. »I've lived on th' moor with 'em so long. I've watched 'em break shell an' come out an' fledge an' learn to fly an' begin to sing, till I think I'm one of 'em. Sometimes I think p'raps I'm a bird, or a fox, or a rabbit, or a squirrel, or even a beetle, an' I don't know it.«

He laughed and came back to the log and began to talk about the flower seeds again. He told her what they looked like when they were flowers; he told her how to plant them, and watch them, and feed and water them.

»See here,« he said suddenly, turning round to look at her. »I'll plant them for thee myself. Where is tha' garden?«

Mary's thin hands clutched each other as they lay on her lap. She did not know what to say, so for a whole minute she said nothing. She had never thought of this. She felt miserable. And she felt as if she went red and then pale.

»Tha's got a bit o' garden, hasn't tha'?« Dickon said.

It was true that she had turned red and then pale. Dickon saw her do it, and as she still said nothing, he began to be puzzled.

»Wouldn't they give thee a bit?« he asked. »Hasn't tha' got any yet?«

She held her hands even tighter and turned her eyes toward him.

»I don't know anything about boys,« she said slowly. »Could you keep a secret, if I told you one? It's a great secret. I don't know what I should do if any one found it out. I believe I should die!« She said the last sentence quite fiercely.

Dickon looked more puzzled than ever and even rubbed his hand over his rough head again, but he answered quite good-humoredly.

»I'm keepin' secrets all th' time,« he said. »If I couldn't keep secrets from th' other lads, secrets about foxes' cubs, an' birds' nests, an' wild things' holes, there'd be naught safe on th' moor. Aye, I can keep secrets.«

Mistress Mary did not mean to put out her hand and clutch his sleeve but she did it.

»I've stolen a garden,« she said very fast. »It isn't mine. It isn't anybody's. Nobody wants it, nobody cares for it, nobody ever goes into it. Perhaps everything is dead in it already; I don't know.«

She began to feel hot and as contrary as she had ever felt in her life.

»I don't care, I don't care! Nobody has any right to take it from me when I care about it and they don't. They're letting it die, all shut in by itself,« she ended passionately, and she threw her arms over her face and burst out crying—poor little Mistress Mary.

Dickon's curious blue eyes grew rounder and rounder.

»Eh-h-h!« he said, drawing his exclamation out slowly, and the way he did it meant both wonder and sympathy.

»I've nothing to do,« said Mary. »Nothing belongs to me. I found it myself and I got into it myself. I was only just like the robin, and they wouldn't take it from the robin.«

»Where is it?« asked Dickon in a dropped voice.

Mistress Mary got up from the log at once. She knew she felt contrary again, and obstinate, and she did not care at all. She was imperious and Indian, and at the same time hot and sorrowful.

»Come with me and I'll show you,« she said.

She led him round the laurel path and to the walk where the ivy grew so thickly. Dickon followed her with a queer, almost pitying, look on his face. He felt as if he were being led to look at some strange bird's nest and must move softly. When she stepped to the wall and lifted the hanging ivy he started. There was a door and Mary pushed it slowly open and they passed in together, and then Mary stood and waved her hand round defiantly.

»It's this,« she said. »It's a secret garden, and I'm the only one in the world who wants it to be alive.«

Dickon looked round and round about it, and round and round again.

»Eh!« he almost whispered, »it is a queer, pretty place! It's like as if a body was in a dream.«